\chapter{The data}
\label{ch:data}

\section{The goal}

Before anything can be predicted, one has to know what the data actually
contains. So the first question of this report is not ``how well can a model
predict the next antenna?'' but the more modest one that has to be answered
first: \emph{what does one line of this dataset represent, and what kind of
claim can it support?}

The answer shaped every experiment that follows. Two of its consequences, that
\textbf{a record is not a movement} and that \textbf{the name of an antenna is
not a place}, are the reason several later results look surprising at first
sight. This chapter establishes both, and ends with the list of the six sets of
people I trained and tested on.

\section{What one line of the data looks like}

The material is fifteen daily files covering 12 to 26 March 2014 in the Karlovy
Vary study area of north-west Bohemia. They are plain text
files\footnote{\textbf{CSV}, \emph{comma-separated values}: the simplest tabular
text format, one row per line with the fields separated by a delimiter. Here the
delimiter is a semicolon.} in which \textbf{one line is one person for one day},
from midnight to midnight. Depending on the day, a file holds between 80{,}217
and 99{,}504 lines; the day used for most of the work reported here, 12 March
2014, holds 85{,}946 people.

A line starts with eight fixed columns, then repeats a pair of values for as
many records as that person produced during the day:

\begin{itemize}
  \item the \textbf{eight fixed columns} are an anonymous number standing for the
        person, their age, their gender, an unexplained five-digit number, a
        behaviour class, a ``home antenna'', an ``activity antenna'' and the
        number of records that follow;
  \item then, over and over, \textbf{the name of an antenna followed by a time},
        the time always coming just after the antenna it belongs to. The time is
        written as a number of seconds since midnight, from 0 to 86{,}399, so
        dividing it by 3{,}600 and dropping what remains gives the hour of the
        day at which the record was written.
\end{itemize}

Only the anonymous number, the record count and the antenna-and-time pairs are
used in this work. The remaining columns carry known aberrations, including ages
of 1078 and negative values, and none of them is needed to predict a next
antenna.

\subsection*{A record is not a movement}

This is the single most important property of the data, and it is easy to miss.

A record is written whenever a phone and a mast exchange a message. That
includes the case where nothing happens at all: a mast asks every phone attached
to it to announce itself on a fixed cycle, typically every four hours on 2G and
every two hours on 3G,\footnote{\textbf{Re-signalling}, or \emph{periodic
location update}: the network asking every attached phone to announce itself on
a fixed cycle, so that it knows where to route an incoming call. It happens
whether or not anyone touches the phone.} so that it knows where to route an
incoming call. A person who spends the whole day in one place therefore still
produces a stream of records, at a rate set by the network rather than by
themselves.

Two consequences follow immediately, and both matter for the rest of the report.

\begin{itemize}
  \item \textbf{How many records a person has is only partly about that person.}
        It also depends on which generation of radio their phone attached to, and
        on how much they used it.
  \item \textbf{A long run of the same antenna name is the normal case}, not an
        anomaly. Any predictor that simply repeats the last name it saw will be
        right very often, entirely for free.
\end{itemize}

\section{The rule I have to beat, and why it moves}

The second point above is what creates the yardstick this whole report is
measured against, so it is worth stating precisely before any model appears.

The \textbf{persistence baseline} is the rule ``the next antenna is the last one
seen''. It contains no model, no training and no settings. It is scored on
exactly the same moments as my model, in exactly the same way, so the two are
directly comparable.

Its value is not a constant, and this is the second thing that is easy to miss.
It is a property of the people being scored. On a day file taken as it comes,
where a large share of people barely move, ``repeat the last antenna'' is right
69.2\,\% of the time. Once the data is restricted to people who actually move
around, the very same rule falls to 45.5\,\%. Nothing about the rule changed;
the population did.

This is why every result in this report is stated as a \textbf{distance}, my
model minus the rule, both measured on the same people and the same moments, and
never as a bare accuracy. Whenever a chapter introduces a new set of people, the
value of the rule on them is stated with it, and the reader should resist
comparing an accuracy in one chapter with an accuracy in another.

\section{The name of an antenna is not a place}

An antenna name such as \cid{BSOCHO1} looks like an opaque location code. It is
not, and this is the point at which the technical name for it can safely be
introduced: in the working files this string is called the
\textbf{\cid{cell\_id}}, and it is used with that name in the rest of the report.

It decomposes into three parts, which Figure~\ref{fig:cellid-anatomy} takes
apart: a letter giving the generation of radio, a middle part naming the
\emph{mast},\footnote{\textbf{Site root}: the middle letters of a
\cid{cell\_id}, naming the physical mast, \cid{KVCHE} in \cid{UKVCHE102}.
\textbf{Sector}: the trailing digits, naming which antenna of that mast served
the phone.} and trailing digits naming which antenna of that mast was in use.
The study area contains \textbf{369 distinct antennas over 113 masts}, split
into 224 antennas on 2G and 145 on 3G. In other words, the same physical mast
appears in the data under three or four different names.

% F2 -- Anatomy of a cell name. A schema: almost no text inside the drawing.
\begin{figure}[!htbp]
\centering
\fitwidth{%
\begin{tikzpicture}[font=\small,
  seg/.style={draw=clrule,minimum height=1.0cm,inner xsep=8pt,align=center},
  tag/.style={align=center,font=\footnotesize,text width=3.0cm},
  ar/.style={-{Stealth[length=4pt]},clrule},
  chip/.style={draw=clrule,fill=white,rounded corners=2pt,inner sep=4pt,
               font=\footnotesize\ttfamily},
]
% ---------------- the name, taken apart ----------------
\node[seg,fill=clmodel!14] (t)                  at (2.2,0) {\ttfamily\LARGE B};
\node[seg,fill=clink!18,right=0pt of t]  (root) {\ttfamily\LARGE SOCHO};
\node[seg,fill=clmove!18,right=0pt of root] (z) {\ttfamily\LARGE 1};

\node[tag,anchor=north] (nt)    at (0.9,-1.05)
  {\textbf{2G or 3G}\\ \cid{B}, \cid{D} = 2G\\ \cid{U}, \cid{V} = 3G};
\node[tag,anchor=north] (nroot) at (4.5,-1.05)
  {\textbf{which mast}\\ \emph{this, and only this,}\\ \emph{is a place}};
\node[tag,anchor=north] (nz)    at (8.1,-1.05)
  {\textbf{which antenna}\\ of that mast\\ (one digit, or three)};

\draw[ar] (t.south)    -- ++(0,-0.25) -| (nt.north);
\draw[ar] (root.south) -- ++(0,-0.25) -| (nroot.north);
\draw[ar] (z.south)    -- ++(0,-0.25) -| (nz.north);

% ---------------- four names, one mast ----------------
\begin{scope}[xshift=11.35cm,yshift=0.65cm]
  \node[chip] (c1) at (0,0)       {BKVCHE1};
  \node[chip] (c2) at (2.15,0)    {BKVCHE4};
  \node[chip] (c3) at (0,-0.72)   {UKVCHE102};
  \node[chip] (c4) at (2.15,-0.72){UKVCHE201};
  \node[draw=clink,fill=clink!12,rounded corners=2pt,inner sep=5pt,
        font=\footnotesize,text width=4.4cm,align=center] (share) at (1.07,-2.05)
        {four names, one mast: \cid{KVCHE}};
  \draw[ar,clink] (c3.south) -- ++(0,-0.30) -| (share.north);
  \draw[ar,clink] (c4.south) -- ++(0,-0.30) -| (share.north);
\end{scope}
\end{tikzpicture}}
\caption{What the name of an antenna contains}
\label{fig:cellid-anatomy}
\figtext{%
\figpara{The name is not a code, it is a description.}{On the left, the antenna
name \cid{BSOCHO1} taken apart. The first letter says which generation of radio
served the phone, 2G or 3G. The middle letters name the mast, and that middle
part is the only piece of the name that corresponds to a location on a map. The
last digits say which antenna of that mast the phone was on, and since a mast
carries three or four antennas pointing in different directions, those digits
change without anybody going anywhere.}

\figpara{The consequence, on the right.}{The same mast in the town of Cheb
appears in the data under four different names, because it carries two
generations of radio and several antennas each. A phone switching from
\cid{BKVCHE1} to \cid{UKVCHE102} has changed name and not changed place.}

\figpara{Why this decides how everything else is measured.}{On the day used for
the analysis of Chapter~\ref{ch:groups}, \textbf{12.4 out of every 100}
back-to-back records are exactly that: a new antenna name, the same mast, and a
person who did not move. Counting those as movement would inflate every mobility
figure in this report by about a quarter and would scramble the four kinds of
user of Chapter~\ref{ch:groups}, whose entire purpose is to separate people who
move from people who do not. So whenever this report measures \emph{movement},
it compares masts. What the model has to \emph{predict}, on the other hand, stays
the full antenna name, because that is the question the operator asked.}}
\end{figure}


The practical consequence is measurable, and Figure~\ref{fig:transition-composition}
measures it. On the day used for the analysis of Chapter~\ref{ch:groups}, out of
every hundred back-to-back records: 43.6 carry the very same antenna name;
44.0 are a genuine change of mast, so the person really went somewhere; and
\textbf{12.4} fall in between, the phone having re-attached to a different
antenna of the \emph{same} mast without the person moving at all.

% F4 -- What a consecutive-event pair actually is, on two datasets.
\begin{figure}[!htbp]
\centering
\fitwidth{%
\begin{tikzpicture}
\begin{axis}[rep,
  xbar stacked, width=\linewidth, height=3.4cm,
  /pgf/bar width=17pt,
  xmin=0, xmax=100,
  xlabel={share of all consecutive-event pairs (\%)},
  ytick={0,1}, yticklabels={\smob, \sraw},
  y tick label style={align=right},
  grid=none, xmajorgrids,
  nodes near coords, point meta=explicit symbolic,
  every node near coord/.append style={font=\scriptsize,text=white,anchor=center},
  legend style={at={(0.5,-0.70)},anchor=north,legend columns=3,draw=clrule},
  legend cell align=left, legend image code/.code={\draw[#1,draw=none] (0cm,-0.08cm) rectangle (0.28cm,0.14cm);},
]
\addplot[fill=clmodel,draw=none] coordinates {(37,0) [37.0\,\%] (43.6,1) [43.6\,\%]};
\addplot[fill=clink,draw=none]   coordinates {(12,0) [12.0\,\%] (12.4,1) [12.4\,\%]};
\addplot[fill=clmove,draw=none]  coordinates {(51,0) [51.0\,\%] (44.0,1) [44.0\,\%]};
\legend{same antenna name, same mast\,--\,other antenna, the person really moved}
\end{axis}
\end{tikzpicture}}
\caption{The three things a pair of back-to-back records can be}
\label{fig:transition-composition}

\end{figure}


Every measurement of \emph{movement} in this report is therefore taken on the
mast, never on the raw antenna name. The \emph{prediction target}, on the other
hand, remains the raw antenna name. That is the quantity the original question
asks for, and coarsening it to the mast would be answering an easier question
than the one that was posed.

\section{From a raw file to training material}

Turning a day file into something a language model can be trained on takes four
steps, drawn in Figure~\ref{fig:data-flow}: keep only the columns that make up a
trajectory, turn each time into an hour of the day, write one ordered sequence
of events per person, and split the people into a training group and a test
group.

% F3 -- From a raw day-file to training sequences, and the disjoint user split.
\begin{figure}[t]
\centering
\fitwidth{%
\begin{tikzpicture}[font=\small,
  stage/.style={draw=clrule,fill=clsoft,rounded corners=2pt,align=center,
                inner sep=5pt,text width=3.15cm,minimum height=1.5cm},
  raw/.style={draw=clrule,fill=white,rounded corners=2pt,align=left,
              inner sep=5pt,font=\scriptsize\ttfamily,text width=14.4cm},
  disk/.style={draw=clmodel,fill=clmodel!10,circle,align=center,
               inner sep=2pt,minimum size=2.05cm,font=\footnotesize},
  disktest/.style={draw=clink,fill=clink!14,circle,align=center,
               inner sep=2pt,minimum size=2.05cm,font=\footnotesize},
  ar/.style={-{Stealth[length=4pt]},clrule,thick},
]
\node[raw] (raw) {892739;22;M;36452;CETR;;;136;BKVDOU3;23822;BKVVIV2;24105;BKVSTN4;24154;\dots};
\node[font=\scriptsize,anchor=north,text=clrule] at ([xshift=-4.9cm,yshift=-0.34cm]raw.south)
  {$\underbrace{\hspace{4.6em}}$};
\node[font=\scriptsize,anchor=north,text=clrule] at ([xshift=-4.9cm,yshift=-0.62cm]raw.south)
  {8 metadata columns};
\node[font=\scriptsize,anchor=north,text=clrule] at ([xshift=1.9cm,yshift=-0.34cm]raw.south)
  {$\underbrace{\hspace{16em}}$};
\node[font=\scriptsize,anchor=north,text=clrule] at ([xshift=1.9cm,yshift=-0.62cm]raw.south)
  {(\cid{cell\_id}, timestamp) pairs --- the timestamp always \emph{follows} its cell};

\node[stage,below left=1.95cm and 3.4cm of raw.south] (s1)
  {\textbf{1. Keep what is a\\trajectory}\\[1pt]\footnotesize user id, record count,\\ the (cell, timestamp) pairs};
\node[stage,right=0.5cm of s1] (s2)
  {\textbf{2. Second $\rightarrow$ hour}\\[1pt]\footnotesize
   $\lfloor \mathit{ts}/3600 \rfloor$\\ one hour token per event};
\node[stage,right=0.5cm of s2] (s3)
  {\textbf{3. One line per user}\\[1pt]\footnotesize an ordered sequence of
   (\cid{cell\_id}, hour) events for one day};
\node[stage,right=0.5cm of s3] (s4)
  {\textbf{4. Split by user}\\[1pt]\footnotesize stratified on the number of
   records, \emph{before} any training};

\draw[ar] ([yshift=-0.95cm]raw.south) -- ++(0,-0.35) -| (s1.north);
\draw[ar] (s1) -- (s2); \draw[ar] (s2) -- (s3); \draw[ar] (s3) -- (s4);

\node[disk,below=1.25cm of s2] (tr) {\textbf{TRAIN}\\[2pt]{\large 400}\\[1pt]users the\\model studies};
\node[disktest,below=1.25cm of s4] (te) {\textbf{TEST}\\[2pt]{\large 100}\\[1pt]users it is\\graded on};
\node[font=\footnotesize,text=clloss,align=center] (neq) at ($(tr)!0.5!(te)$)
  {$\neq$\\[1pt]\textbf{0 user}\\ in common};

\draw[ar] (s4.south) -- ++(0,-0.4) -| (te.north);
\draw[ar] (s4.south) -- ++(0,-0.4) -| (tr.north);

\node[draw=clrule,fill=white,rounded corners=2pt,align=left,inner sep=6pt,
      text width=14.4cm,font=\footnotesize,below=0.75cm of neq]
  {\textbf{What the split guarantees, and what it does not.} The model never
   sees a single test user before being graded, so nothing it scores is
   memorised. But a user is \emph{one day}, and the day-files share no users at
   all: the model can only ever learn how a \emph{population} moves, never how
   one person moves. This single fact explains most of the ceilings reported in
   Chapters~\ref{ch:user} and~\ref{ch:scale}.};
\end{tikzpicture}}
\caption[From a raw day-file to training sequences]{How a raw daily export
becomes training material. The 400/100 figures are those of the
mobile-user day; the other tiers of Table~\ref{tab:datasets}
follow the same four stages with different counts.}
\label{fig:data-flow}
\end{figure}


Two choices inside that pipeline deserve to be stated, because they change how
later results should be read.

\paragraph{The split is by person, and it is balanced.} Training and test people
are different people: the model is graded only on strangers. The assignment is
not made at random but is \emph{balanced} on the number of records per
person,\footnote{\textbf{Stratified split}: people are first sorted into bands by
how many records they have, and the split is made inside each band, so both
sides end up with the same mix of short and long histories.} and once the split
is made the two sides are compared to confirm that they really are consistent
with each other. Without that precaution, a split can silently hand all the
easier, longer histories to one side.

\paragraph{Cleaning removes people, not records.} Three successive filtering
decisions were made over the internship, each one in response to a result rather
than in advance.

\begin{itemize}
  \item \textbf{Nothing at first.} The earliest runs used the day file exactly as
        it came. This turned out to be why the rule sat at 69.2\,\%: a large
        share of people produce a handful of records from a single antenna, so
        the evaluation was dominated by moments where there was nothing to
        predict.
  \item \textbf{More than ten records, and no four-hour hole.} The four-hour
        threshold is the corpus's own convention: since a mast re-signals at
        least every four hours, a gap longer than that means the person was
        outside the recording area rather than sitting still inside it.
  \item \textbf{People who move, and who stay in the area}, using the per-day
        selection prepared by a colleague. This is the data on which most of the
        modelling rests, and it is where the rule drops to the mid-forties.
\end{itemize}

Each of these makes the \emph{task} harder and the \emph{rule} weaker, which is
exactly the point. A benchmark on which one line of code already scores 69\,\%
leaves almost nothing to measure.

\section{The six sets of people I used}

Table~\ref{tab:datasets} lists every set of people used over the eleven
iterations, with the counts measured on the files themselves and the score the
one-line rule reaches on each. Six sets for one problem may look like
indiscipline; it is the direct consequence of the point above, that each new
idea needed data the earlier files did not carry. Adding the time of day
required regenerating the sequences with their times kept. Testing whether more
people help required a day file with more people in it. Testing whether
\emph{many} more people help required stacking eleven day files end to
end.\footnote{\textbf{Stacking} day files: putting them one after another into a
single training file, which adds \emph{people} rather than adding days to the
people already there.} Each regeneration produced a new test group, and
therefore a new value for the rule.

% T1 -- Every set of people used, measured on the files themselves.
\begin{table}[!htbp]
\centering\footnotesize
\caption{The six sets of people, and the score of the one-line rule on each}
\label{tab:datasets}
\begin{tabular}{@{}l@{\hspace{5pt}}>{\raggedright\arraybackslash}p{4.6cm}@{\hspace{5pt}}r@{\hspace{6pt}}r@{\hspace{6pt}}l@{\hspace{5pt}}l@{}}
\toprule
& \textbf{Who is in it} & \textbf{Train} & \textbf{Test} & \textbf{Records/person} & \textbf{Rule} \\
\midrule
\sraw  & one day file taken as it comes (12 March 2014) & 400 & 100 & mean $\approx$62 & 69.2\,\% \\
\sfilt  & the same day, filtered: more than 10 records and no 4-hour hole & 400 & 100 & $>$10 by design & 79.0\,\% \\
\smob  & people who move, and who stay in the area, on one day & 400 & 100 & 19--200, mean 86.0 & 45.5 / 43.1\,\% \\
\slev & the same people, every day stretched to exactly 200 records & 400 & 100 & 200 exactly & 76.5\,\% \\
\stwok  & a day file sampled down to 2{,}000 people & 900 / 1{,}900 & 100 & 12--200, mean 77.1 & 48.2\,\% \\
\seleven  & 11 day files of 2{,}000 people, stacked end to end & 21{,}000 & 1{,}000 & 10--200 & 50.2\,\% \\
\bottomrule
\end{tabular}

\vspace{4pt}
\vspace{4pt}
\vspace{4pt}
\end{table}


\section{What the data cannot answer}

Two limits follow from the data itself, and neither can be fixed by better
modelling.

\paragraph{The unit is one day.} The question posed at the start of this
internship is a within-day one: predict the rest of a person's day from its
beginning. The data is built for exactly that, one person is one day, and every
experiment reported here answers that question. Predicting tomorrow from the
days before is a different question, and the natural continuation of this work;
Chapter~\ref{ch:conclusion} returns to it.

\paragraph{The world is small.} 369 antennas over 113 masts, with a typical
person visiting on the order of twenty distinct antennas. Already at the
smallest training size, 98.9\,\% of the antennas appearing in the test people's
days had also appeared in the training people's, and from 5{,}000 people onward
the figure is 100\,\%. So ``discovering an antenna it had never seen'' cannot be
a source of improvement beyond that point, and any remaining gain has to come
from better estimates of the moves between antennas that are already known.

\section{Looking back}
\label{sec:retro-data}

The decision that mattered most in this chapter was not a modelling decision at
all. It was accepting that the raw day file could not support the question being
asked. Moving to the selection of mobile, non-absent people that Nino extracted
for me was the first real lever of the internship, and I would describe it as a
data-quality result rather than a modelling one: it produced trajectories on
which \emph{analysing the past in order to predict the rest} is a meaningful
exercise at all. It cost the comfortable-looking accuracies of the early runs,
and that is precisely what made every chapter after it measurable.

The one thing I would sequence differently is the order of those two steps. I
spent the first iterations improving the training machinery on data that could
not answer the question, and only then changed the data. Settling that first
would have reached the same conclusions sooner.
